%%%%%%%%%%  scan idx 54  |  folio 41  %%%%%%%%%%
\begin{equation}\tag{2.8}
 C^1(L(D))=\ellc(D)\qquad C^i(L(D))=0\quad (i\geq 1)
\end{equation}
(où
pour un cycle quelconque $Z$, on note $\ellc(Z)$ sa classe dans $A(X)$). Enfin, on impose que les $C^i$ soient \emph{fonctoriels}, i.e. que si $f$ est un morphisme de $X$ dans $Y$ et $E$ un fibré vectoriel sur $Y$, on ait
\begin{equation}\tag{2.9}
 C^i(f^!(E))=f^*(C^i(E)),
\end{equation}
où $f^!(E)$ est l'image inverse de $E$ sur $X$. On peut aussi écrire (2.9) sous la forme suivante : la notion d'image inverse de fibré définit un $\lambda$-homomorphisme de $\lambda$-anneaux augmentés
\begin{equation}\tag{2.10}
 f^!:\underline{K}(Y)\longrightarrow \underline{K}(X)\footnote{La notation $f^!$ utilisée ici est en conflit avec celle utilisée en théorie de la dualité (cf par exemple R. Hartshorne, Residues and Duality, Lecture Notes in Math. n\textsuperscript{o} 20, Springer, 1966). C'est pourquoi dans ce Séminaire nous la remplaçons par la notation $f^*$.}
\end{equation}
(on écrit ici $\underline{K}(X)$ au lieu de $\underline{K}(X)$), et (2.9) signifie qu'on a commutativité dans le diagramme~:
\[
\begin{tikzcd}[column sep=large,row sep=large]
K(Y) \arrow[r,"f^!"] \arrow[d,"\Ctil_Y"'] & K(X) \arrow[d,"\Ctil_X"]\\
A(Y) \arrow[r,"f^*"] & A(X).
\end{tikzcd}
\]
(on a mis en indice au symbole $\Ctil$ l'espace sur lequel on le considère.)

\textbf{Théorème 2.2.} \emph{Il existe une théorie des classes de Chern satisfaisant aux conditions (2.5), (2.7), (2.8) et (2.9).}

Nous admettrons aussi la proposition suivante~:

\textbf{Proposition 2.3.} \emph{Soit $Y$ une partie fermée de $X$ et soit $U$ son complémentaire~; si $x$ est un élément de $A(X)$ qui induit $0$ sur $U$, il existe un représentant de $x$ qui soit un cycle porté par $Y$.}

